\documentclass[12pt,twoside]{book}

\usepackage[utf8]{inputenc}
\usepackage[turkish]{babel}
\usepackage[T1]{fontenc}

\usepackage{amsmath}
\usepackage{amssymb}
\usepackage{amsthm}
\usepackage{enumerate}
\usepackage[most]{tcolorbox}

\usepackage{circuitikz}
\usepackage{cancel}
\usepackage{bodegraph}
\usepackage{gensymb}

\usepackage{verbatim}
\usepackage{listings}
\usepackage{matlab-prettifier}

\usepackage{pgfplots}
\pgfplotsset{compat=newest}
\usepackage{pgffor}
\usepackage{xcolor}

\usepackage{tikz}
\usetikzlibrary{patterns}
\usetikzlibrary{babel}

\usetikzlibrary{shapes.misc}
\tikzset{cross/.style={cross out, draw=black, minimum size=2*(#1-\pgflinewidth), inner sep=0pt, outer sep=0pt},
%default radius will be 3pt. 
cross/.default={3pt}}
\usetikzlibrary{arrows}
\tikzset{
  arrow/.pic={\path[tips,every arrow/.try,->,>=#1] (0,0) -- +(.1pt,0);},
  pics/arrow/.default={triangle 90}
}

\usepackage{siunitx}
\usepackage{textcomp}

\usepackage{hyperref}
\hypersetup{
    colorlinks=true,
    linkcolor=blue,
    filecolor=magenta,      
    urlcolor=cyan,
    pdftitle={Kitap},
    pdfpagemode=FullScreen,
    }

\urlstyle{same}

\usepackage{graphicx}
\graphicspath{{svg-inkscape/}}
\usepackage{geometry}
\usepackage{subcaption}
\usepackage{tabularray}
\usepackage{tabularx}
\DefTblrTemplate{contfoot-text}{default}{Bir sonraki sayfada devam ediyor.} % Uzun Tablolar için alt kısım
\DefTblrTemplate{conthead-text}{default}{(Devam)} % Uzun Tablolar için üst kısım

\usepackage{lscape}

\usepackage{pgfplotstable}
\usepackage{booktabs}

\renewcommand{\chaptername}{Bölüm}
\renewcommand{\contentsname}{İçindekiler}
\renewcommand{\figurename}{Şekil}
\renewcommand{\tablename}{Tablo}
\renewcommand{\bibname}{Kaynaklar}
\renewcommand{\listfigurename}{Şekiller}
\renewcommand{\listtablename}{Tablolar}
\renewcommand{\appendixname}{Ek}
\renewcommand{\indexname}{Dizin}
\renewcommand{\partname}{Bölüm}
\renewcommand{\proofname}{İspat}

\begin{document}
\shorthandoff{=} % babel paketi hatasını gidermek için gerekli

\tableofcontents

\documentclass[17pt]{beamer}

\usepackage[utf8]{inputenc}
\usepackage[turkish]{babel}
\usepackage[T1]{fontenc}

\usepackage{amsmath}
\usepackage{amssymb}
\usepackage{amsthm}
\usepackage{enumerate}
\usepackage{tikz}
\usepackage{transparent}
\usepackage{xcolor}

\usetheme{Warsaw}
\usecolortheme{default}
\definecolor{nigdeyesili_acik}{RGB}{3, 150, 166}
\definecolor{nigdeyesili_koyu}{RGB}{143, 209, 217}
\setbeamercolor{structure}{fg=nigdeyesili_acik}

\author[Dr. Mehmet CANEVİ]{Arş.~Gör.~Dr.~M.~Canevi\inst{1}}

\institute{
    \inst{1}%
    Bilgisayar Mühendisliği\\
    Mühendislik Fakültesi
}
    
\date[2025] {Ders Notları, Ocak 2025}
    
\logo{\transparent{0.4}\includegraphics[width = 20mm]{logo}}
    
\title[Ders 1] {Kod yazma, Kod derleme. Değişken tanımlama, veri tipleri. Operatörler.}
\begin{document}
%%%%%%%%%%%%%%%%%%%%%%%%%%%%%%%%%%%%%%%%%%%%%%%%%%%%%%%%%%%%%%%%%%%%%%%%%%%%%%%%
\frame{\titlepage}
\begin{frame}[fragile]{İçidekiler}
    \tableofcontents
\end{frame}
%%%%%%%%%%%%%%%%%%%%%%%%%%%%%%%%%%%%%%%%%%%%%%%%%%%%%%%%%%%%%%%%%%%%%%%%%%%%%%%%
\section{Örnek}
\begin{frame}[fragile]
    \lstinputlisting{lec1/main.c}
\end{frame}
%%%%%%%%%%%%%%%%%%%%%%%%%%%%%%%%%%%%%%%%%%%%%%%%%%%%%%%%%%%%%%%%%%%%%%%%%%%%%%%%
\section{Kod yazma}
\begin{frame}[fragile]{Kod nedir?}
    Bilgisayar tanımı kapsamına giren cihazlara meramımızı anlatmak için yazılan şeye \textbf{kod} diyebiliriz.
\end{frame}
%%%%%%%%%%%%%%%%%%%%%%%%%%%%%%%%%%%%%%%%%%%%%%%%%%%%%%%%%%%%%%%%%%%%%%%%%%%%%%%%
\section{Değişken tanımlama ve veri tipleri}
\begin{frame}[fragile]{Değişken isimleri}
    Değişkenler şu şekilde tanımlanmaktadır:
    \begin{lstlisting}
        veri_tipi degisken_adi;
        veri_tipi degisken_adi=ilk_deger;\end{lstlisting}
    Buradaki \lstinline{veri_tipi};
    \begin{itemize}
        \item \lstinline{int} 
        \item \lstinline{float}
        \item \lstinline{double}
        \item \lstinline{char}
    \end{itemize}
     olabilmektedir. \lstinline{degisken_adi} ise 
    \begin{itemize}
     \item türkçe karakter içermemelidir
     \item sadece harf veya alttire ile başlamalıdır
     \item özel karakterler(!,?,vb.) ve boşluk içermemelidir
     \item \lstinline{veri_tipi} anahtar kelimeleri olmamalıdır
    \end{itemize}
    kurallarına uymalıdır.
\end{frame}
%%%%%%%%%%%%%%%%%%%%%%%%%%%%%%%%%%%%%%%%%%%%%%%%%%%%%%%%%%%%%%%%%%%%%%%%%%%%%%%%
\begin{frame}[fragile]{Değişken tiplerinin anlamı}
    Değişkenler
    \begin{itemize}
     \item \lstinline{int} ingilizce \verb|integer| yani tam sayıdan gelmektedir \lstinline{-10,4} gibi
     \item \lstinline{float} ingilizcede kayan noktalı sayı(6-7 basamak) anlamındadır \lstinline{3.14f} gibi
     \item \lstinline{double} ingilizcede iki kat hassasiyet(15 basamak) anlamındadır \lstinline{3.14} gibi
     \item \lstinline{char} ingilizce \verb|character|'den gelir \lstinline{'c','x'} gibi
    \end{itemize}
    veri tipleri ile tanımlanabilmektedir ve 
    \begin{itemize}
        \item \lstinline{int var1=-10;} 
        \item \lstinline{float var2=3.14f;} 
        \item \lstinline{double var3=3.14;} 
        \item \lstinline{char var4='x';} 
    \end{itemize}
    olarak örneklendirilebilir.
\end{frame}
%%%%%%%%%%%%%%%%%%%%%%%%%%%%%%%%%%%%%%%%%%%%%%%%%%%%%%%%%%%%%%%%%%%%%%%%%%%%%%%%
\begin{frame}[fragile]{Çoklu değişken tanımlama}
    Birden fazla değişken
    \begin{itemize}
        \item \lstinline{int a,b,c=-10;} 
        \item \lstinline{float a=2.2,b,c=3.14f;} 
        \item \lstinline{double a,b,c=3.14;} 
        \item \lstinline{char a,b,c;a=b=c='x';}
    \end{itemize}
    şeklinde tanımlanabilir.
\end{frame}
%%%%%%%%%%%%%%%%%%%%%%%%%%%%%%%%%%%%%%%%%%%%%%%%%%%%%%%%%%%%%%%%%%%%%%%%%%%%%%%%
\begin{frame}[fragile]{Değişkenlerin yazdırılması}
    Değişkenler
    \begin{itemize}
        \item \lstinline{int var1=-10;printf("%d",var1);} 
        \item \lstinline{float var2=3.14f;printf("%f",var2);} 
        \item \lstinline{double var3=3.14;printf("%lf",var3);} 
        \item \lstinline{char var4='x';printf("%c",var4);}
        \item \lstinline{printf("%s","Bu bir cumledir.");}
    \end{itemize}
    ile ekrana yazdırılır.
\end{frame}
%%%%%%%%%%%%%%%%%%%%%%%%%%%%%%%%%%%%%%%%%%%%%%%%%%%%%%%%%%%%%%%%%%%%%%%%%%%%%%%%
\begin{frame}[fragile]{Sabit değerli değişkenler}
    Değeri değişmemesi gereken değişkenler
    \begin{itemize}
        \item \lstinline{const int SAAT_DK=60;} 
        \item \lstinline{const float PI_SAYISI=3.14;} 
    \end{itemize}
    şeklinde tanımlanır ve büyük harflerin kullanımı bir gelenektir. Dolayısıyla,
    \begin{itemize}
        \item \lstinline{SAAT_DK=120;} 
        \item \lstinline{PI_SAYISI=5.5;} 
    \end{itemize}
    hata verir.
\end{frame}
%%%%%%%%%%%%%%%%%%%%%%%%%%%%%%%%%%%%%%%%%%%%%%%%%%%%%%%%%%%%%%%%%%%%%%%%%%%%%%%%
\section{Operatörler}
\begin{frame}[fragile]{Matematiksel operatörler}
    Toplama ve çıkarma işlemi
    \begin{itemize}
        \item \lstinline{int a=4+5;int c=a-2;int d=a+c;} 
    \end{itemize}
    ile örneklendirilebilir. Matematiksel operatörler
    \begin{itemize}
        \item \lstinline{+,-} toplama, çıkarma 
        \item \lstinline{*,/} çarpma, bölme
        \item \lstinline{++,--} arttırma, azaltma 
        \item \lstinline{%} \verb|mod| operatörü (bölümden kalan)
        \item \lstinline{=} atama operatörü
    \end{itemize}
    şeklindedir.
\end{frame}
%%%%%%%%%%%%%%%%%%%%%%%%%%%%%%%%%%%%%%%%%%%%%%%%%%%%%%%%%%%%%%%%%%%%%%%%%%%%%%%%
\begin{frame}[fragile]{Matematiksel operatörler(devam)}
    Aşağıdaki operatörler tanımlıdır.
    \begin{itemize}
        \item \lstinline{x+=3;} \lstinline{x=x+3;} 
        \item \lstinline{x-=3;} \lstinline{x=x-3;} 
        \item \lstinline{x*=3;} \lstinline{x=x*3;} 
        \item \lstinline{x/=3;} \lstinline{x=x/3;} 
        \item \lstinline{x%=3;} \lstinline{x=x%3;} 
    \end{itemize}
\end{frame}
%%%%%%%%%%%%%%%%%%%%%%%%%%%%%%%%%%%%%%%%%%%%%%%%%%%%%%%%%%%%%%%%%%%%%%%%%%%%%%%%
\begin{frame}[fragile]{Karşılaştırma operatörleri}
    Aşağıdaki operatörler ile karşılaştırma başarılı ise 1 değilse sıfır değerini alır.
    \begin{itemize}
        \item \lstinline{x>y} Büyüktür
        \item \lstinline{x>=y} Büyük eşittir
        \item \lstinline{x<y} Küçüktür
        \item \lstinline{x<=y} Küçük eşittir
        \item \lstinline{x==y} Eşittir
        \item \lstinline{x!=y} Eşit değildir
    \end{itemize}
\end{frame}
%%%%%%%%%%%%%%%%%%%%%%%%%%%%%%%%%%%%%%%%%%%%%%%%%%%%%%%%%%%%%%%%%%%%%%%%%%%%%%%%
\begin{frame}[fragile]{Mantıksal operatörleri}
    Mantıksal operatörler
    \begin{itemize}
        \item \lstinline{sart1 && sart2} VE
        \item \lstinline{sart1 || sart2} VEYA 
        \item \lstinline{!sart} DEĞİLDİR 
    \end{itemize}
    olarak verilmiştir.
\end{frame}
%%%%%%%%%%%%%%%%%%%%%%%%%%%%%%%%%%%%%%%%%%%%%%%%%%%%%%%%%%%%%%%%%%%%%%%%%%%%%%%%
\end{document}
\chapter{Ayrıklaştırma}
Türevin geometrik yorumu 
\begin{equation}
    \frac{dy(t)}{dt}\approx\frac{\Delta y}{\Delta t}
\end{equation}
olmak üzere
\begin{equation}
\begin{split}
    \frac{dy(t)}{dt}&\approx\frac{\Delta y}{\Delta t}\\
    &\approx\frac{y((k+1)T)-y(kT)}{(k+1)T-kT}\\
    &\approx\frac{y((k+1)T)-y(kT)}{T}
\end{split}
\end{equation}
elde edilir. Ayrık bir sinyalin türevi ardışık değerler farkının örnekleme zamanına oranı ile hesaplanabilmektedir. Örneğin, $y(kT)=\sin(kT)$ ve $T=0.1$ olmak üzere
\begin{equation}
    \frac{y((k+1)T)-y(kT)}{T}=10(\sin((k+1)0.1)-\sin(0.1k))
\end{equation}
ve dolayısıyla
\begin{equation}
\begin{split}
    \{10\sin(0.1),10(\sin(0.2)-\sin(0.1)),10(\sin(0.3)-\sin(0.2)),\cdots\}\\
    \{0.9983,0.9884, 0.9685,\cdots\}
\end{split}
\end{equation}
elde edilir. $y(kT)=\sin(kT)$ sinyalinin türevinin $\frac{d\sin(t)}{dt}=\cos(t)$ olduğu bilindiğinden
\begin{equation}
    \begin{split}
        \{cos(0.1),cos(0.2),cos(0.3),\cdots\}\\
        \{ 0.9950,0.9801,0.9553,\cdots\}
    \end{split}
\end{equation}
elde edilir ve ayrık türev ile benzer değerler olduğu görülmektedir. Bu yaklaşıklığın türeve yakınsaması için örnekleme zamanı $T$ daha küçük seçilmelidir. 
\begin{equation}
    \frac{dq(t)}{dt}=x
\end{equation}
olmak üzere
\begin{equation}
\begin{split}
    \frac{dq(t)}{dt}&=x\\
    dq(t)&=xdt\\
    \int dq(t)&=\int xdt\\
    q(t)&=\int xdt
\end{split}
\end{equation}
elde edilir. Buradan hareketle,
\begin{equation}
    \begin{split}
        \frac{\Delta q}{\Delta t}&=x\\
        \frac{q((k+1)T)-q(kT)}{(k+1)T-kT}&=x\\
        \frac{q((k+1)T)-q(kT)}{T}&=x\\
        q((k+1)T)-q(kT)&=xT\\
        q((k+1)T)&=q(kT)+xT
    \end{split}
\end{equation}
ifadesi bulunur. Ayrık zamanda integral birikimli toplama karşılık gelmektedir. Bu karşılıklar Zero Order Hold(ZOH) ile elde edilmiştir. ZOH örnekleme zamanı boyunca değerlerin sabit olduğu varsayımına dayanmaktadır. Bu durum
\begin{equation}
    x(t)=x(kT),\quad kT\leq t\leq (k+1)T
\end{equation}
ile ifade edilebilir.

First Order Hold(FOH) yöntemi ise
\begin{equation}
    x(t)=x(kT)+\frac{t-kT}{T}(x((k+1)T)-x(kT)),\quad kT\leq t\leq (k+1)T
\end{equation}
olarak tanımlanır. Eşitliğin sağ tarafı $t=kT$ için $x(kT)$, $t=(k+0.5)T$ için 
\begin{equation}
\begin{split}
    x(t)&=x(kT)+\frac{t-kT}{T}(x((k+1)T)-x(kT)),\quad kT\leq t\leq (k+1)T\\
    &=x(kT)+\frac{kT+0.5T-kT}{T}(x((k+1)T)-x(kT))\\
    &=x(kT)+0.5(x((k+1)T)-x(kT))\\
    &=x(kT)+0.5x((k+1)T)-0.5x(kT)\\
    &=0.5x((k+1)T)+0.5x(kT)
\end{split}
\end{equation}
ve $t=(k+1)T$ için ise 
\begin{equation}
    \begin{split}
        x(t)&=x(kT)+\frac{t-kT}{T}(x((k+1)T)-x(kT)),\quad kT\leq t\leq (k+1)T\\
        x(t)&=x(kT)+\frac{(k+1)T-kT}{T}(x((k+1)T)-x(kT))\\
        x(t)&=x(kT)+x((k+1)T)-x(kT)\\
        x(t)&=x((k+1)T)
    \end{split}
\end{equation}
elde edilir. Görüldüğü üzere ZOH yönteminin aksine $T$ süre boyunca değerler değişmektedir.

\end{document}
